% before: th-ranks
The auxiliary lattice of the argument is built from a chain of $m$ theta lifts
between definite unitary groups, and every later estimate depends on a handful
of integers. The first of them is the Galois datum of
Chapter~\ref{chapter-host-fields}, on which the host field is built; the others
fix the size of the starting group, the number of lifts and the rank added by
each lift. The inequalities among them are the ones the later chapters need,
for instance that each block is longer than the number of lifts.

% before: th-hc-parameter
The forms that are lifted have prescribed archimedean types. These are recorded
by half-integral and integral sequences, one for the starting group and one
pair of integers for each lift. Interlacing of the first two sequences is the
condition under which the relevant Gan-Gross-Prasad period can be nonzero.

% before: th-spaces
The groups of the tower are unitary groups of Hermitian spaces over the host
field. Each space carries a basis, a signature and a unimodular lattice, since
the later arguments count lattice classes and compute archimedean
representations. The definite spaces $X_j$ come with indefinite companions
$Y_j$, which receive them isometrically and on which the indefinite side of
the argument works.

% before: th-column-chain
An algebraic modular form on a definite unitary group is a function on the
finite set of classes of lattices in its genus. The lift from one space of the
tower to the next is a finite sum over pairs of classes, weighted by the
inverse orders of the stabilizers. The following data are what such a sum
needs, together with exponents of determinants that shift the archimedean
weight by the right amount at each step.

% before: th-iterated-lift
With the data of a column chain the lifts can be written down. Each lift
multiplies by a power of the determinants $\delta_i$, applies the kernel of a
pair of classes, and divides by another power of the determinants. Composing
$m$ of them gives the map whose periods are measured in the rest of the
argument.

% before: th-definite-datum
The argument works with one structure that holds the field, the spaces, a
column chain for each word, the periods of the iterated lifts and the further
data used to measure them. Collecting these in one place lets every later
statement quantify over a single kind of object. The weights in part (7) make
visible the block shape of the iterated lift.

% before: th-twisting-data
At the places of the coefficient field where the tuple of Hecke characters
differs from the base tuple, the argument of Chapter~\ref{chapter-assembly}
twists the forms by characters of ray class groups. The datum carries, for each
such twist, the Galois representations, theta elements and theta values whose
behaviour under the twist is controlled there. Their lattices, Bloch-Kato
lengths and theta values are defined here so that the twist condition can be
stated.

% before: th-lift-kernels
Nothing in the datum itself says that the kernels are theta kernels. That
identification matters only where the integrality of the lifted forms is
proved. There each kernel is required to be a differential operator applied at
the origin to a combination of theta functions of the pair of lattices.

% before: th-weil-datum
Theta lifts between unitary groups are defined through the Weil
representation. The model used here consists of functions on the whole adelic
symplectic space that are quasi-periodic under its rational points, so that
the theta kernel is evaluation at the origin. The Weil representation is
specified by its properties, which is what the seesaw arguments below use.

% before: th-grid-places
The theta grid defined below carries Hecke characters that may ramify only at
a few places of $F$: a base place above an auxiliary prime $q_0$ and the places
above the level primes. At each of these places the grid chooses a place of $E$
that is not conjugate to itself, so that the local idempotents exist. The
choices must agree with the places of the level data and with the ramification
of the characters of the level system.

% before: th-grid-centres
At the archimedean places the characters of a theta grid are fixed by integers
built from the Harish-Chandra parameter: the midpoints between consecutive
terms of an arithmetic progression above $v_1(\tau)$. The step of the
progression is large compared with the ranks of the grid, which makes the
archimedean exponents of the characters of each square of the grid large
compared with the ranks of its spaces.

% before: th-grid
The norms of the lifts are computed on the indefinite side, where the definite
spaces sit inside a grid of Hermitian spaces of growing rank. Moving along a
row of the grid lifts forms by one block of the tower; moving along a column
lifts by a fixed amount that does not depend on the tuple of characters. The
level system and the grid carry the Hecke characters that enter the $L$-values
below.

% before: th-bottom-triple
The first lifts in the grid are lifts of the trivial representation from
totally definite groups. At the bottom of each square of the grid these take
place between three spaces with identity Gram matrices, and the theta
functions are fixed at the places where the characters ramify. Their theta
integrals are the objects whose norms and nonvanishing are needed.

% before: th-large-data
The statements about bottom triples hold only once the scale is large, and the
largeness is uniform in the remaining data. The same quantifier pattern occurs
several times, so it is stated once. The $L$-functions that appear are partial
Hecke $L$-functions, assumed entire.

% before: th-rallis-inner-product
The covolume bound of Chapter~\ref{chapter-covolumes} needs the norms of the
first theta lifts. The Rallis inner product formula expresses such a norm as a
product of Hecke $L$-values at consecutive integers. It is used here for the
lift of the trivial representation of a definite group.

% before: th-theta-hypotheses
The nonvanishing of the first lifts uses two global facts about theta integrals
beyond the published local inputs. They compare theta integrals of the two
halves of a square and track how the comparison changes along a Witt tower.
Both are stated in the form in which they enter.

% before: th-first-lift-matching
The second half of the nonvanishing theorem compares the theta integrals of
the two halves of a bottom triple. Both halves map onto one space of
coinvariants of the Weil representation, on which the unitary group of the
large space acts irreducibly. A matching records this space with the two maps
onto it, and a square triple records the conditions on ranks, signatures and
characters under which matchings exist.

% before: th-square-fock-operators
On the vectors $f\otimes P$ with a fixed archimedean vacuum $f$, the Lie
algebras of the two unitary groups act through the polynomial $P$. They act by
quadratic expressions in raising and lowering operators, shifted by the
archimedean exponents of the characters. Howe duality at the real places is a
finiteness statement about the subspaces that are stable under one Lie algebra
and contain the coinvariants of the other, and the radical is the sum of the
proper ones.

% before: th-square-fock-derivatives
A Weil datum differentiates along the one-parameter subgroups at a complex
place into the operators just defined. The finite adelic points of the larger
group act on a vector $f\otimes P$ through its vacuum $f$ alone, since they
commute with the archimedean Heisenberg action.

% before: th-square-fock-maps
With the vacuum fixed, a map of a matching becomes a linear map on the
polynomial space. The conditions of a matching make it equivariant for the
Lie algebra of the larger group and make it vanish on the coinvariants of the
smaller group, and the finite adelic points carry one such map to another.

% before: th-square-howe-model
To apply Howe duality at the real places, the polynomial space has to be
identified with the Fock space of the published theorem. Under the
identification, subspaces stable under the Lie algebra of the larger group
become stable under its Lie algebra and maximal compact subgroup there, and
the coinvariants contain the image of the action of the smaller group.

% before: th-square-howe-condition
The identification turns subspaces of the polynomial space that are stable and
contain the coinvariants into theta subspaces. Since it is injective and onto,
the two conclusions of Howe duality at the real places become the two halves of
the Howe condition.

% before: th-square-radical
The archimedean factor of a matching is the quotient of the polynomial space by
the radical. Four facts of linear algebra produce it: the radical is the
largest proper admissible subspace, a space spanned by images of intertwiners
and irreducible under a larger family is semisimple, the radical lies in the
kernel of every intertwiner, and two such quotients that map onto the same
irreducible module are isomorphic.

% before: th-darboux-model
The local matching and the agreement of theta integrals are statements about
Schwartz-Bruhat functions on a Lagrangian half of the symplectic space, in
Darboux coordinates. A function there gives a theta function on the whole
space by summing its Heisenberg translates over the rational points, and the
algebraic functions are the pure tensors with a Fock vector at the
archimedean places and the vacuum at almost all finite places.

% before: th-darboux-reading
A square triple and its Weil data have to be read as a triple over $E/F$
before the hypotheses on theta integrals apply to them. The base data come
from the ramification of the characters, the spaces keep their Gram matrices
up to sign and transposition, and the Weil representations of the model are
transported from the Weil data along the theta-sum map.

% before: th-kernel-theta-subspace
The kernel of a matching and the kernel of the matched local maps are both
subspaces of the algebraic functions that are stable under the larger group
and contain the coinvariants of the smaller one. The images of algebraic
functions under the matched local maps are formed factor by factor, over a
finite set of places outside which the functions are the vacua.

% before: th-kernel-dictionary
The conditions of a matching concern the finite adelic points and the
derivatives along one-parameter subgroups at the complex places. Through the
theta-sum map they become conditions on the algebraic functions of the model,
for the compact subgroups and one-parameter subgroups of the unitary groups
of the model, and the conditions on matched local maps become conditions on
the images of expressions.

% before: th-kernel-matching
Since the space of a matching has no proper stable subspaces, the kernel of a
matching, read on algebraic functions, cannot be enlarged to another proper
theta subspace. This maximality is what lets Howe duality identify it with the
kernel of the matched local maps.

% before: th-howe-finite
At the real places Howe duality is Theorem~\ref{lit-howe-duality-real}. At the
finite places the comparison of kernels needs a local statement about the
theta lift of the trivial representation of the smaller group: either it is
irreducible, or it is generated by one vector whose stabilizer averages every
vector into the line of that vector.

% before: th-howe-radical
For the comparison of kernels, each finite place enters through the kernel of
its matched local map. This kernel is the largest proper stable subspace, and
the local form of Howe duality gives a generating vector and averaging
operators that fix it and carry finitely many vectors of the kernel into the
coinvariants.

% before: th-howe-archimedean
At the real places the factor is the archimedean Fock space, and Howe duality
at the real places bounds the chains of its theta subspaces. Together with the
irreducibility of the space of the matched maps there, it puts every proper
theta subspace into the kernel of the real matched map.

% before: th-howe-tensor
Algebraic functions supported at a finite set of places form a tensor product
of the archimedean Fock space with local Schwartz-Bruhat spaces. A subspace of
it that is stable under the factorwise operators and contains the factorwise
relations is either everything or lies in the kernel of the tensor product of
the local maps; finite length gives this at the real places, and the radical
condition of the finite places replaces finite length there.

% before: th-kernel-howe
The tensor product statement applies to every finite set of places at once.
A proper theta subspace that contained a function with nonzero image would
contain all algebraic functions supported at every larger finite set of
places, which is all of them.

% before: th-kernel-graph
Once the kernels agree, the pairs of images of functions with agreeing images
under the matched local maps form a subspace of the square of the space of a
matching that is the graph of an intertwining map. Countable dimension and
irreducibility make this map a scalar.

% before: th-assemble-matching
Matched local maps for a triple give a space of coinvariants: the tensor
product of their spaces at the real places with the restricted tensor product
of the local spaces. The theta-sum map identifies the algebraic functions of
the model with the functions of a matching, and each condition of a matching
is a condition on matched local maps read factor by factor.

% before: th-square-model
The local matching and the agreement of theta integrals concern Alpöge's model
of a triple over $E/F$, built on Schwartz-Bruhat functions in Darboux
coordinates. A square triple and its Weil data are read in this model through
the theta-sum map. The lemma lists what the reading gives: the Weil data, the
algebraic functions, the theta integrals, and the comparison of a matching
with matched local maps.

% before: th-bottom-square-facts
The bottom triples of the grid are definite and carry the standard archimedean
data. Their ranks, the archimedean exponents of their characters and the
ramification of these characters are fixed by the construction of the grid and
the level system. These facts make a bottom triple a square triple and
describe its places of ramification.

% before: th-bottom-split-reading
To compare $f_t$ and $f'_t$ through matched local maps, the bottom triple is
read in Alpöge's model with its own lattices, so that its Fock vacuum vectors
are theta sums of pure tensors of the local vacua of the model. The embeddings at the
places below $S_t$ are chosen to define the places of $S_t$. The finite set of
places of the resulting triple is then the set of places of $F$ below $S_t$.

% before: th-bottom-split-assembly
The matching assembled from matched local maps is computed on theta sums of
pure tensors, whose images are tensor products of local images. The
computation uses of the triple only the properties of
Lemma~\ref{th-square-model}(2), so it applies to the triple read with the
lattices of the bottom triple.

% before: th-bottom-split-zeta
At a place below $S_t$ the matched local maps are compared through zeta
functions on general linear groups. The local components of $f_t$ and $f'_t$
there are standard functions built from the conductor of the character of the
square. Their zeta functions agree after a change of rows and columns and a
nonzero scalar, by a local computation.

% before: th-bottom-split-situation
The next two lemmas concern one bottom triple, its reading in Alpöge's model
and local maps matched to it. At the places below $S_t$ they use linear models
of the Schrödinger models, in which half of the Heisenberg group acts by
translations and the other half by characters. The conditions on matched local
maps that the argument uses are listed with the setup.

% before: th-bottom-split-local
At a place below $S_t$, the zeta matching of the local maps turns the local
computation of Lemma~\ref{th-bottom-split-zeta} into an equality of local
images. The changes of rows and columns compare the split bases of the model
with the frames of the bottom triple, so the local functions obtained are the
components of $f_t$ and $f'_t$ at that place, read in linear models.

% before: th-bottom-split-pure
The vectors $f_t\otimes P$ are built from a Fock vacuum vector by replacing
one finite component and applying creation operators, and the images under the
assembled matching are known on theta sums of pure tensors. The lemma writes
$f_t\otimes P$ and $f'_t\otimes P'$ as such theta sums, with the local
functions of Lemma~\ref{th-bottom-split-local} as components below $S_t$. The
argument extends the vacuum place by place, using the irreducibility of the
Schrödinger representation, and then applies the creation operators.

% before: th-bottom-split-matching
At the places of $S_t$ the functions $f_t$ and $f'_t$ have standard
components built from the conductor of the character of the square. Reading
the bottom triple in Alpöge's model with these components gives a matching in
which either every image of a Fock vector of $f_t$ vanishes, or one of them is
a multiple of an image of a Fock vector of $f'_t$.

% before: th-first-lift-inputs
The two hypotheses on theta integrals are stated for Alpöge's model of a
triple over $E/F$, which does not refer to the spaces of the grid. Carrying
them to square triples, and computing the bottom triples at the split places
where their characters ramify, gives the following six facts. The rest of the
argument for the second lift uses them together with linear algebra and
integration over compact quotients.

% before: th-first-lift-dixmier
The archimedean factor of the coinvariants has countable dimension and is
irreducible under the Lie algebra alone. By Dixmier's form of Schur's lemma
its commutant is then the scalars, and this is what splits the coinvariants
as a tensor product.

% before: th-first-lift-product
A space spanned by the images of intertwiners from an irreducible module with
scalar commutant is the tensor product of that module with the space of
intertwiners. Applied to both halves of a square triple, this gives one
multiplicity space for the two halves.

% before: th-first-lift-factor
The archimedean factor and the product form together give the factorization
(F1) used in the nonvanishing theorem. The archimedean factor depends only on
the spaces and the archimedean exponents, so for a bottom triple it is the
same for every tuple. A matching with given Weil data comes from any matching,
since Weil data with the same characters agree on theta functions.

% before: th-first-lift-compact
At an archimedean place the compact subgroup acts on the Fock vectors through
a linear change of the variables of the polynomials, twisted by powers of
determinants. In particular it does not raise the degree.

% before: th-first-lift-relations
For a definite group the theta integrals converge, and they vanish on the
relations by which coinvariants for the finite adelic group are formed. The
finiteness of the archimedean orbits for a definite target group is the third
analytic fact that the irreducibility argument needs.

% before: th-first-lift-comparison
The agreement of theta integrals compares the two halves of a square triple
only after the second function is changed by a relation. The relations lie in
the kernel of the coinvariant map and the theta integral vanishes on them, so
the comparison holds for all pairs of functions with the same image.

% before: th-first-lift-comparison-short
The proof above follows Alpöge's order, in which the vanishing for functions
with zero image is proved on its own. The comparison itself already contains
it, because the theta integral of the zero function is zero.

% before: th-first-lift-irreducible
The functions whose theta integral vanishes map to a subspace of the
coinvariants that is stable under the finite adelic group and the Lie
algebra. Irreducibility leaves two cases, which give the dichotomy (F3).

% before: th-first-lift-unique
Two matchings with the same Weil data have the same kernels, so on each half
they differ by an isomorphism of the spaces of coinvariants. The two
isomorphisms differ by an operator that commutes with an irreducible action,
hence by a scalar.

% before: th-first-lift-finite
At the finite places the test functions of a bottom triple are explicit, and
for one matching the images of the two standard vectors are proportional. The
uniqueness of matchings carries this to every matching and every
factorization, which gives (F4).

% before: th-first-lift-nonzero
A period of the iterated lift can only be nonzero if the lifts themselves are
nonzero. For the lifts of the trivial representation at the bottom of the grid
this follows from the inner product formula, once the $L$-values in it are
known not to vanish, by the criterion of Rallis and of Kudla and Rallis
\cite{Ral84,KR94}. The statement covers both halves of every square.

% before: th-tensor-rule
The seesaw arguments split theta functions along orthogonal decompositions of
either space of a pair. A theta function for an orthogonal sum can be built
from theta functions for the summands, and the Weil representation of the sum
acts factor by factor. For Weil data given by their properties this needs a
proof: the two actions can differ only by a character, which the conditions on
Weil data force to be trivial.

% before: th-line-comparison
The seesaw for one step of the tower matches a line in one space of a pair with
a line in the other space of another pair. A linear map between the lines with
the right norm identifies the corresponding symplectic spaces up to sign. Under
this identification theta functions and the actions of the Weil data on them
correspond, with a complex conjugation when the sign requires it.

% before: th-seesaw-identity
The periods of the iterated lift are compared with Bessel periods through
seesaws, one for each step of the tower. The basic identity is an equality of
products of theta kernels for two pairs of spaces that share a line. It holds
pointwise, for arbitrary Weil data whose characters fit together as stated.

% before: th-matrix-coefficients
The orthogonality relation below integrates the coefficients
$\langle f,\rho(w)f_0\rangle$ of two theta functions over the whole of
$\mathbb W$. That integral converges because the coefficients are supported, in
the finite directions, near a lattice, and decay rapidly in the archimedean
directions. Summing them over the rational translates of a point then gives a
bound that does not depend on the point.

% before: th-lattice-vacua
Vanishing statements for functionals on theta functions are proved by expanding
a theta function in translates of a fixed one. The expansion is an
orthogonality relation for the translations of the Heisenberg group. On lattice
vacua it reduces to the translations by the annihilator of the lattice, since
the other translates are orthogonal to all lattice vacua.

% before: th-vacuum-irreducibility
The seesaw arguments end by showing that a linear functional on theta
functions vanishes identically. The space of lattice vacua is irreducible for
the translations by the annihilator of the lattice, so it suffices to check the
vanishing on the translates of one nonzero vector. Tameness supplies the
continuity and the interchange of integrals that this needs.

% before: th-heights
The growth of Weil operators is measured by heights on the adelic unitary
groups. The height of an element bounds its entries and those of its inverse at
every place, and the height on the quotient takes the best rational translate.
Sizes of rational vectors seen from a group element give a second measure,
which reduction theory compares with the height.

% before: th-archimedean-growth
At the archimedean places a Weil datum is the oscillator representation up to a
scalar of absolute value one. In the Fock model the oscillator representation
is an integral of translates against the vacuum projection, whose size is a
determinant. This gives polynomial growth of theta functions under elements
with finite component one.

% before: th-sized-lattice
The reduction of the next lemmas moves the lattice of an adelic element to a
fixed lattice by a rational isometry. This needs an orthogonal basis of the
lattice with norms in a fixed finite set and sizes bounded by a power of the
size of a spanning set. The hypotheses on the lattice used in the construction
are collected here.

% before: th-anisotropic-frame
The first step finds an orthogonal basis with nonzero norms and polynomially
bounded sizes, by Gram-Schmidt within the lattice. Nothing yet controls the
norms themselves, which may be large.

% before: th-short-vector
To control the norms, a Minkowski argument on the trace form finds a nonzero
lattice vector whose norm has bounded conjugates. The orthogonal basis of the
previous lemma makes the trace form diagonal over the real subfield, which is
what bounds the size of this vector.

% before: th-anisotropic-vector
A short vector may be isotropic. In that case it is completed to a hyperbolic
pair inside the lattice, and a vector of the pair adjusted by a multiple of the
other has its norm in a fixed finite set of residues.

% before: th-orthogonal-frame
Passing to orthogonal complements, the anisotropic vectors of the previous
lemma build an orthogonal basis with norms in a fixed finite set. The
complement of such a vector inherits the hypotheses with a controlled change of
constants.

% before: th-reduction
To pass from the archimedean places to the whole adelic group, every element is
moved by a rational element into a fixed finite set of cosets of an open
subgroup, times an archimedean element. The archimedean part must be bounded
polynomially in the height, so that the growth of the previous lemma survives.

% before: th-height-comparison
The tameness of the periods along the tower requires the heights of embedded
elements to be comparable with the heights in the smaller groups. Both
directions are proved through the sizes of rational vectors: a lower bound for
these sizes bounds the height on the quotient, and embedding the vectors
transports such bounds.

% before: th-weil-estimates
The functionals to which the irreducibility lemma is applied are periods of
theta kernels, integrated against rapidly decreasing functions. Checking that
they are tame needs four properties of the Weil operators on theta functions:
how they move translations, joint continuity, polynomial growth in the
heights, and commutation with integrals of translates. Heights must also be
compared along the embeddings of the decompositions used in the tower.

% before: th-kernel-tame
Every period in the rest of the chapter is an integral, over a space of
parameters, of values of Weil operators applied to a theta function. One
criterion covers all of them: the weight must decay faster than any power of
the heights of the group elements. The tameness of each period then reduces to
an estimate of its weight.

% before: th-seesaw-period-nonvanishing
The ascent of nonvanishing starts at the bottom of the tower. There a period
of a theta lift restricted to a smaller unitary group is turned into a
Fourier-Jacobi period of depth zero. The proof combines the seesaw with the
irreducibility of lattice vacua.

% before: th-periods
Above the bottom step the periods that occur are Bessel and Fourier-Jacobi
periods, the two kinds of Gan-Gross-Prasad periods. They integrate a pair of
forms over a smaller unitary group and a unipotent group, against a character
or a theta function.

% before: th-core-chain
The ascent concerns two families of forms on the grid. One is a chain of forms
on the core spaces, each obtained from the next by a theta lift. The other is
the sequence of lifts of one form along the bottom row.

% before: th-step
Above the bottom of the tower each step compares a Bessel period on one space
with a Fourier-Jacobi period on the next, or the reverse. The comparison runs
through a Weil datum for the pair of spaces and one for the pair of their
anisotropic parts. A theta function for the large pair is carried to one for the
small pair by partial Fourier transforms along the hyperbolic parts followed by
an average over a unipotent group.

% before: th-step-local
The comparison of the two periods is an unfolding computation, which needs a
supply of theta functions closed under the right translations. These are local
conditions at the finite places: invariance under a compact open subgroup, a
support condition, and a lattice for which the transformed functions are
vacua.

% before: th-step-periods
Unfolding a Bessel period of a theta lift does not produce a Fourier-Jacobi
period directly. It produces a functional of the transformed theta function in
which the lift from the smaller unitary group is still unevaluated. The two
functionals below play this role for the two kinds of steps.

% before: th-step-functionals
The vanishing arguments apply Lemma~\ref{th-vacuum-irreducibility} to the
period functionals, which requires them to be tame. On products of theta
functions the auxiliary functionals reduce to the periods of the next step, by
the tensor rule. The transformed theta functions also commute with the
Heisenberg translations of the small pair.

% before: th-mixed-model
The unfolding expands a theta function along the isotropic part of a
hyperbolic decomposition. In this mixed model the elements of the parabolic
subgroup act by explicit formulas: the Levi factor by a change of variables
with a character, the unipotent elements by Heisenberg translations and a
central character. The terms of the theta series are the values of the mixed
model at rational points.

% before: th-mixed-model-literature
The formulas of the mixed model are familiar from the local theory of the Weil
representation. Here they hold for any Weil datum, because the proof uses only
the conditions that define one.

% before: th-mixed-origin
Moving the point of the mixed model to the origin reduces the formulas to the
value of the mixed model at $0$. There the operators of the parabolic subgroup
and of the unitary group of $W$ act through the Weil datum of the smaller pair,
up to characters. The characters are pinned down at the rational points and at
the split places.

% before: th-unfolding
The steps of the tower compare a period of a theta lift with a period of the
flat maps of the theta functions. The comparison is an unfolding: the
integral over the flag group picks out, in the expansion of the theta kernel,
the terms indexed by one orbit of rational points, and that orbit unfolds the
integral over $U(W)(F)\backslash U(W)(\mathbf A)$. The same computation serves
the Bessel and the Fourier-Jacobi periods, so it is stated for a general weight
on the flag group.

% before: th-bessel-unfolding
The unfolding lemma needs a weight on the flag group with three properties,
and for the Bessel period this weight is the conjugate of the Bessel
character. Its central values are trivial because the base point has Gram
matrix zero, and transvections provide the nontrivial values for the dependent
terms. The unfolded functional is then the flat map evaluated at zero.

% before: th-fourier-strips
The vanishing of a Bessel period for all translates of a theta function by
archimedean elements gives the vanishing of the unfolded integrand against all
characters of the archimedean coordinates. Fourier uniqueness turns this into
vanishing on sets defined by those coordinates, and strips of the Jacobi group
are such sets.

% before: th-vanishing-descends
Each step rests on an unfolding computation. If the period of the lift
vanishes for all theta functions of the local data, the unfolded integrals
vanish on strips, and a continuity argument gives the vanishing of the
auxiliary functional at the transformed theta functions.

% before: th-step-nonvanishing
Combining the unfolding with the irreducibility of lattice vacua turns the
vanishing on the transformed theta functions into vanishing on all lattice
vacua. The product formulas then identify one of these values with the period
of the next step. Read contrapositively, this carries nonvanishing from one
step to the next.

% before: th-grid-conjugation
The theta grid is set up with the first space of each pair negated, while the
steps of the tower use the spaces themselves. Complex conjugation passes from
one convention to the other. It inverts the characters, and it turns the
column chain and the partial row lifts into theta lifts for the conjugated Weil
data.

% before: th-lattice-signs
Translations by elements of a self-dual lattice act on its lattice vacua. The
places of the base field over two split in the host field, so every lattice
element is a sum of elements in split halves and of twice a lattice element,
and on lattice vacua each such translation acts by a sign.

% before: th-coset-vectors
The row data are built from coset vectors, and the ascent needs to know how
translations act on them. Coset vectors are lattice vacua for the depth
lattice, and a translation by an element of its annihilator moves a coset
vector to a multiple of another one. These two facts let a functional that
vanishes on coset vectors vanish on all lattice vacua.

% before: th-invariant-form
At a split place, linear forms invariant under one half of the symplectic space
are the only information the Weil datum gives. To show that they detect a
theta function, one must construct such forms that do not vanish on a given
function, and this lemma does so by a positivity argument.

% before: th-split-separation
At a split place the Weil datum is known only through its action on
functionals invariant under one half of the symplectic space. A theta function
on which all these functionals vanish, after translations by the other half,
is zero, so this knowledge determines the action on theta functions.

% before: th-nonsplit-vacua
At a nonsplit place the axioms of a Weil datum do not describe the action of
the lattice stabilizers directly. A Weil datum built from an unramified
character has the unramified property, and any Weil datum with characters
unramified at the place differs from a twist of it by characters that are
trivial on the relevant determinants.

% before: th-depth-fixed
The local data of a step require the function lifted along a row to be
invariant under a compact open subgroup. The theta functions of the row data
are vacua for the depth lattices of the grid, and such vacua are fixed by the
congruence subgroups. So each partial row lift is right invariant under them.

% before: th-fock-translates
A functional that vanishes on the polynomials in the creation operators
applied to a Fock vacuum vector should vanish on its archimedean translates.
The proof is the reproducing property of the Fock model, in which the
translates of the vacuum are Gaussians times entire functions whose Taylor
coefficients are these polynomials. Tameness is what allows the functional to
be moved inside the reproducing integral.

% before: th-fock-span
The steps produce some lattice vacuum for which the next period is nonzero,
while the row data allow only Fock vectors built on coset vectors. A tame
functional that vanishes on all of these vanishes on every lattice vacuum. So a
nonzero value on a lattice vacuum gives a nonzero value on some row datum.

% before: th-tower-decomposition
To run the steps along the row of the grid, the spaces $V_{i,0}$ must be
decomposed compatibly with the core spaces. These decompositions are chosen
once for all positions in the tower. Their conditions at the two ends of the tower
and the local data of every step are stated together.

% before: th-standard-pieces
The local data of a step are built from five groups of coordinates: two in the
hyperbolic part of each space and one in the anisotropic parts. Translations by
these groups act on the mixed models by characters, so conditions imposed on a
theta function through them pass to its flat map. The unfolding conditions
record what the nonvanishing of the flat map forces on an element of the second
unitary group.

% before: th-standard-class
Admissible pieces give the conditions of local data that concern single theta
functions. Each of them reduces to the commutation of translations or to a
computation in the two mixed models. The nonvanishing of the flat map comes from
a theta function built in the mixed models, whose flat map is an integral with
positive real part.

% before: th-standard-local-data
Once the pieces are admissible and split the lattice, the remaining conditions
of local data concern the compact group $K$. They follow from the invariance of
lattice vacua under $K$ and from a statement on the cosets of $R_x$ and $K$.

% before: th-standard-pieces-grid
For the tower decomposition of Lemma~\ref{th-standard-tower} the pieces are cut
out by the depth lattice of the row pair and by a second lattice, in which the
lattice of the second space is scaled by $\mathfrak m\bar{\mathfrak m}^{-1}$. That
ideal has norm $1$, so the second lattice is self-dual, and the pieces along the
hyperbolic part of the second space are each other's annihilators.

% before: th-standard-one-coset
The condition of local data on the support of the flat map is a statement about
cosets in an adelic unitary group. It is proved place by place. Outside the
distinguished places it follows from the Iwasawa decomposition and the
self-duality of the lattices; at the distinguished places, which split in $E$,
it is a computation in a general linear group.

% before: th-standard-lattices
The last conditions of a tower decomposition are statements about lattices. The
anisotropic piece of each step splits along the decomposition of $V'_\rho$ into a
core space and a line, which reduces its vacua to vacua for the depth lattice of
a core pair and for $\Gamma_L(\rho)$.

% before: th-standard-tower
The induction along the tower needs one tower decomposition that fits both
ends and carries local data at every step. It is built from the bases of the
complements of the core spaces given by the grid, scaled by powers of
$-\xi_0$ so that the vectors $\zeta_\rho$ are linked as
Definition~\ref{th-tower-decomposition} requires.

% before: th-tower-rungs
The ascent alternates between Fourier-Jacobi and Bessel periods as the position
in the tower increases by one. Each alternation is one step of
Proposition~\ref{th-step-nonvanishing}, applied to the column chain and to the
row lifts, followed by a choice of row datum. The first step uses the seesaw of
Theorem~\ref{th-seesaw-period-nonvanishing} instead.

% before: th-bessel-ascent
Combining the steps gives the statement used in the argument. A nonzero period
at the bottom of the tower produces a nonzero Bessel period at the top, for the
$m$-fold lift along the bottom row. In the timeline of Alpöge's announcement
it belongs to the seesaw step, in which the period of the iterated lift is
related to a Bessel period.

% before: th-bessel-identity
The division argument of Chapter~\ref{chapter-integral-structure} needs more
than nonvanishing. It needs the period of the iterated lift to be a fixed
multiple of a standard Bessel period that does not depend on the tuple of Hecke
characters, and its values on the division lattices to be rational. These two
requirements are the following condition on a datum.

% before: th-scope
The four parts of this chapter are established with different generality.
Some statements hold for all Hermitian spaces and Weil data over the host
field, others only for the data supplied in Chapter~\ref{chapter-assembly}.
One part of the expected picture, the global parameter of the iterated lift,
is not stated at all.
